\documentclass[11pt]{article}
\usepackage[margin=1in]{geometry}
\usepackage[T1]{fontenc}
\usepackage[utf8]{inputenc}
\usepackage{lmodern}
\usepackage{microtype}
\usepackage{booktabs}
\usepackage{longtable}
\usepackage{amsmath}
\usepackage{xcolor}
\usepackage{enumitem}
\usepackage[numbers,sort]{natbib}
\usepackage[hidelinks]{hyperref}
\usepackage{url}

% Values not yet measured. Every \tbd must be gone before submission.
\newcommand{\tbd}{\textcolor{red}{\textbf{TBD}}}
\newcommand{\todo}[1]{\textcolor{red}{[TODO: #1]}}
\newcommand{\cov}{\texttt{covered}}
\newcommand{\ncov}{\texttt{not\_covered}}
\newcommand{\insuf}{\texttt{insufficient\_documentation}}

\title{MedCoverBench: Can Language Models Decide Medicare Coverage\\from the Policy Text?}
\author{Manoj Dobbali\\Independent researcher\\\url{https://manojdobbali.com}}
\date{Draft, \today}

\begin{document}
\maketitle

\begin{abstract}
Medicare decides whether to pay for a service by applying written coverage policy to the facts of a
case, and the program has begun testing artificial intelligence in that review. We introduce
MedCoverBench, an open benchmark that asks a language model to make this decision from the real
text of a CMS National Coverage Determination (NCD): decide whether a requested service is
covered, not covered, or lacks required documentation; quote the policy language that decides
it; and list what is missing. The benchmark has 513 synthetic cases over 50 NCDs, each label
derived from an explicit checklist of 434 criteria anchored to verbatim policy spans. We report
the \emph{unsafe approval rate}, the share of non-coverable cases a model approves, as the
headline metric, since a ranking by accuracy hides it. Across eight models, the four strongest
reach 91--97\% accuracy with at most 2.7\% unsafe approvals, but small open models approve up to
46\% of cases the policy denies, and a 72B model with 85\% accuracy still approves 18\%. A
long-context track hides the governing NCD in a library of 45 to 140 other policies (about 32k
and 96k tokens), where the model must find the policy before applying it. There, errors shift
toward approval: Llama 3.3 70B still names the governing NCD in 93--98\% of cases, but its unsafe
approval rate rises from 8.7\% with the policy alone to 27.5\% at 32k tokens and 50\% at 96k. All labels are drafts pending review by clinicians and
certified coders. Data are released under CC BY 4.0 and code under Apache 2.0.
\end{abstract}

\section{Introduction}

A large share of healthcare administration is reading a policy and deciding whether a case meets
it. For Original Medicare, the most authoritative of these policies are National Coverage
Determinations (NCDs): documents published by the Centers for Medicare \& Medicaid Services (CMS)
that state when an item or service is reasonable and necessary, for whom, under what conditions
and with what documentation. In 2026 CMS began a six-year model, WISeR, in which technology
vendors, using artificial intelligence and machine learning alongside clinical review, assess
prior authorization requests for selected services in six states~\citep{cms2025wiser}. Whether
language models can apply coverage policy correctly is therefore no longer hypothetical.

Existing benchmarks do not measure this directly. Medical question answering
benchmarks~\citep{jin2021medqa,arora2025healthbench,warner2026medmarks} test clinical knowledge
and conversation, not whether a model follows a specific written policy. Recent agent benchmarks
for healthcare operations~\citep{chen2026chibench,bedi2026healthadminbench} test long workflows
across simulated applications, so a failure can come from navigation or tool use as easily as
from misreading the rule, and the most policy-dense of them is grounded in a synthetic handbook
rather than real coverage policy~\citep{chen2026chibench}. Policy-following benchmarks outside
healthcare~\citep{yao2024taubench,guha2023legalbench,koreeda2021contractnli} show that following
written rules is a distinct and unsolved capability, but none uses coverage policy.

MedCoverBench isolates the decision. Each item gives a model the full text of one real NCD and a
synthetic case, and asks for one of three decisions, the policy quotes that decide it, and any
missing documentation. The third label, \insuf{}, matters in practice: a reviewer who cannot
tell ``the patient does not qualify'' from ``the record does not show that the patient
qualifies'' will deny claims that should be returned for documentation.

Our contributions are:
\begin{itemize}[leftmargin=*,itemsep=2pt]
  \item \textbf{A benchmark built from real coverage policy.} 513 cases over 50 current NCDs,
    with labels derived from 434 checklist criteria, each tied to a span copied verbatim from
    the NCD and checked automatically (\S\ref{sec:data}).
  \item \textbf{A safety-oriented metric suite.} Unsafe approval rate, over-denial rate,
    citation recall and macro-F1 over three decisions (\S\ref{sec:task}).
  \item \textbf{A long-context track} that hides the governing NCD among dozens of other real
    NCDs at controlled length and depth, separating finding the policy from applying it
    (\S\ref{sec:long}).
  \item \textbf{Baselines for eight models} from 7B open models to frontier APIs, showing that
    unsafe approval separates models that accuracy ranks closely (\S\ref{sec:results}).
  \item \textbf{An open release} of data, harness, raw model outputs and a label-review tool,
    with the labels' draft status stated plainly (\S\ref{sec:limits}).
\end{itemize}

\section{Related work}

\paragraph{Medical benchmarks.} MedQA~\citep{jin2021medqa} and its successors test clinical
knowledge through exam questions; HealthBench~\citep{arora2025healthbench} grades open-ended
health conversations against physician-written rubrics; Medmarks~\citep{warner2026medmarks}
aggregates 30 medical benchmarks into an open evaluation suite. These measure what a model knows
about medicine, not whether it applies a payer's written rule. MedCoverBench deliberately forbids
clinical knowledge the policy does not state.

\paragraph{Healthcare administration and agents.} $\chi$-Bench~\citep{chen2026chibench} evaluates
agents on prior authorization, utilization management and care management workflows grounded in a
1,279-document managed-care handbook; the handbook is synthetic. HealthAdminBench~\citep{bedi2026healthadminbench}
evaluates computer-use agents on 135 administrative tasks, including prior authorization, in
simulated EHR and payer portals. Both report low end-to-end success. MedCoverBench is
complementary: it removes the workflow and tests the coverage decision alone, on real CMS text,
so errors can be attributed to policy reasoning.

\paragraph{Following written rules.} $\tau$-bench~\citep{yao2024taubench} shows that agents given
domain policies in context still violate them, and that removing the policies lowers performance.
LegalBench~\citep{guha2023legalbench}, ContractNLI~\citep{koreeda2021contractnli} and
CUAD~\citep{hendrycks2021cuad} test reasoning over legal and contractual text. Coverage policy
shares their structure (conditions, exceptions, documentation requirements) and adds a safety
asymmetry: wrongly approving and wrongly denying have different costs.

\paragraph{Long context and attribution.} Retrieval accuracy in long inputs depends on where the
relevant text sits~\citep{liu2024lost}, and synthetic needle-in-a-haystack tests overstate usable
context length~\citep{hsieh2024ruler,yen2025helmet}. LongBench~\citep{bai2024longbench} and
HELMET~\citep{yen2025helmet} evaluate realistic long-context tasks. Our long-context track is a
realistic variant: the ``haystack'' is other coverage policies, many from the same clinical area,
and the task after retrieval is a rule application, not extraction. Our citation metric follows
the spirit of ALCE~\citep{gao2023alce}, checking that quoted policy text supports the decision.

\section{Data}
\label{sec:data}

\subsection{Source}
NCD text was downloaded from the CMS Coverage API~\citep{cmscoverageapi} on 2026-10-27: 345
current NCDs. A model sees three fields of each NCD: the item or service description, the
indications and limitations of coverage, and reasons for denial. NCD text is a work of the US
government. We exclude Local Coverage Determinations and billing articles because they contain
CPT content licensed by the American Medical Association, which cannot be redistributed; no CPT
codes appear in the cases.

\paragraph{Selection.} We ranked NCDs by how many explicit, checkable conditions their text
contains (``only if'', ``must'', numbered criteria) and excluded policies that are blanket
non-coverage, left to contractor discretion, or covered only within an approved study, since
those cannot support several distinct decisions. Fifty NCDs were selected, spanning 16 chapters
of the NCD manual; cardiovascular services (chapter 20) are the largest group with 14 policies.
Their text is 600 to 3,700 tokens long (median about 1,400).\footnote{Token counts in this paper
are estimated as characters divided by four unless stated otherwise.}

\subsection{Checklists}
For each NCD we extract a checklist of atomic criteria. Each criterion has an identifier, a
paraphrase, a type, and a span copied verbatim from the NCD. The 434 criteria are mostly
eligibility conditions (209), exclusions (69), documentation requirements (61) and provider
qualifications (47), with smaller numbers of device, frequency, exception, setting and
continuation criteria. Where the policy text is ambiguous, the ambiguity is recorded in a
reviewer note and no case may depend on it.

\subsection{Cases}
Cases are written by controlled perturbation of the checklist, so each label follows from it:
\begin{itemize}[leftmargin=*,itemsep=1pt]
  \item every applicable criterion met and documented $\Rightarrow$ \cov{};
  \item one criterion contradicted by the facts $\Rightarrow$ \ncov{}, with that criterion as the
    gold citation;
  \item one criterion's evidence absent, nothing contradicted $\Rightarrow$ \insuf{}, with the
    missing criterion as the gold citation.
\end{itemize}
Authors were instructed to make cases discriminating: values exactly at a threshold that still
qualify, a violation buried among many satisfied details, exceptions the policy grants, and
details that sound disqualifying but are not under this NCD. Each case names the requested
service, which is narrower than the policy's scope where needed (for example ``manual hospital
bed'' rather than the NCD's ``hospital beds, including variable height and electric
features''). Each case has a short rationale.
Patients are synthetic; no protected health information is involved. Vignettes have a median of
85 words (29--251).

The final set has 513 cases: 174 \cov{}, 249 \ncov{} and 90 \insuf{}, with 8 to 13 cases per
policy. Sixteen cases cite more than one criterion.

\paragraph{Automatic validation.} A validator rejects any policy whose checklist or cases fail
three checks: every span appears verbatim in the NCD text (normalizing only quotes, whitespace and
hyphens); every cited criterion exists; and each label is consistent with its citations (a
\cov{} case cites none, a \ncov{} case cites at least one violated criterion, an \insuf{} case
cites only missing ones). Every released policy passes.

\paragraph{Provenance.} Checklists and cases were drafted with AI assistance (Claude, in Claude
Code) under the rules above, in parallel batches of about ten policies, and checked by the
validator. We state this because it bears on how much to trust the labels before expert review
(\S\ref{sec:limits}). During construction we found that our HTML cleaning dropped text between
literal ``<'' and ``>'' characters (for example ``LVEF <35\%''); after the fix, four NCDs gained
text, and we confirmed that no label changed. We report it as a caution for anyone building on
CMS API text.

\section{Task and metrics}
\label{sec:task}

The model receives a fixed system prompt defining the three decisions, the NCD text, the requested
service and the vignette. It must return JSON with a \texttt{decision}, a list of
\texttt{policy\_quotes}, a list of \texttt{missing\_documentation} and a short rationale
(Appendix~\ref{app:prompt}). We take the last valid JSON object in the reply, so reasoning
models may think before answering; a reply with no valid decision is counted as wrong.

\begin{itemize}[leftmargin=*,itemsep=2pt]
  \item \textbf{Unsafe approval rate} (headline): of the cases whose gold label is \ncov{} or
    \insuf{}, the share predicted \cov{}.
  \item \textbf{Over-denial rate}: of the \cov{} cases, the share predicted anything else.
  \item \textbf{Accuracy} and \textbf{macro-F1} over the three decisions.
  \item \textbf{Citation recall}: of the gold criteria in non-covered cases, the share matched by
    at least one quote, where a quote matches if it contains at least half of the criterion
    span's word tokens.
  \item \textbf{Parse failure rate}.
\end{itemize}

We treat unsafe approval as the headline because its costs are asymmetric in practice: an
over-denial is usually appealable, while an unsafe approval pays for a service the policy does not
cover and may expose a patient to an unnecessary procedure.

\section{Long-context track}
\label{sec:long}

In practice a reviewer is not handed the one relevant policy. The long-context track gives the
model a library of NCDs, does not say which applies, and asks for the same output plus the number
of the NCD applied.

\paragraph{Construction.} We sample 120 cases, 40 per decision, spread across all 50 policies, and
exclude seven cases that every strong baseline missed, since they are the likeliest label errors
(\S\ref{sec:results}). For each case we build a library from the other current NCDs: about half from
the target's chapter (for a defibrillator case, other cardiovascular policies) and the rest from
other chapters. We exclude the target's parent and child NCDs (e.g.\ 20.9 and 20.9.1) and any
NCD that cross-references the target or that the target cross-references, since these can govern
the same service. The library is filled to about 32k or 96k tokens; the 32k library for a case
is a subset of its 96k library, so a change in score between sizes reflects length rather than a
different set of distractors. The target's position cycles through relative depths 0, 0.25, 0.5,
0.75 and 1. Libraries hold 45 NCDs on average at 32k and 139 at 96k.

\paragraph{Metric.} In addition to the metrics above we report \textbf{NCD accuracy}: the share of
replies that name the governing NCD. A model that names the wrong NCD but decides correctly is
counted correct on the decision, so NCD accuracy and decision accuracy can be read separately.

\section{Experiments}
\label{sec:results}

\paragraph{Setup.} We evaluate eight models through OpenRouter: GPT-5 mini, Gemini 3.1 Flash Lite,
Qwen 3.7 Plus, DeepSeek V3.1, Llama 3.3 70B Instruct, Qwen 2.5 72B Instruct, Qwen 2.5 7B Instruct
and Llama 3.1 8B Instruct. All runs use temperature 0 and an 8,000-token output limit, so reasoning
models are not truncated. Each model sees each case once. The full single-policy evaluation of all
eight models cost about US\$4.50, and the long-context runs reported here about \tbd.

\begin{table}[t]
\centering
\small
\caption{Single-policy track, all 513 cases. Unsafe approval is the headline metric (lower is
better). Labels are drafts; see \S\ref{sec:limits}.}
\label{tab:main}
\begin{tabular}{lcccccc}
\toprule
Model & Acc. & Macro-F1 & Unsafe appr.\ $\downarrow$ & Over-denial $\downarrow$ & Cite rec. & Parse fail \\
\midrule
GPT-5 mini            & 0.967 & 0.966 & 0.009 & 0.023 & 0.992 & 0.008 \\
Qwen 3.7 Plus         & 0.943 & 0.927 & 0.009 & 0.057 & 0.963 & 0.000 \\
Gemini 3.1 Flash Lite & 0.934 & 0.930 & 0.015 & 0.115 & 0.927 & 0.000 \\
DeepSeek V3.1         & 0.910 & 0.922 & 0.027 & 0.172 & 0.932 & 0.012 \\
Llama 3.3 70B         & 0.893 & 0.908 & 0.086 & 0.109 & 0.870 & 0.014 \\
Qwen 2.5 72B          & 0.848 & 0.871 & 0.180 & 0.075 & 0.930 & 0.000 \\
Qwen 2.5 7B           & 0.622 & 0.532 & 0.224 & 0.276 & 0.718 & 0.006 \\
Llama 3.1 8B          & 0.591 & 0.571 & 0.463 & 0.092 & 0.670 & 0.016 \\
\bottomrule
\end{tabular}
\end{table}

\begin{table}[t]
\centering
\small
\caption{Recall per gold decision, single-policy track.}
\label{tab:recall}
\begin{tabular}{lccc}
\toprule
Model & \cov{} & \ncov{} & \insuf{} \\
\midrule
GPT-5 mini            & 0.98 & 0.98 & 0.90 \\
Qwen 3.7 Plus         & 0.94 & 0.98 & 0.84 \\
Gemini 3.1 Flash Lite & 0.89 & 0.98 & 0.91 \\
DeepSeek V3.1         & 0.83 & 0.96 & 0.94 \\
Llama 3.3 70B         & 0.89 & 0.89 & 0.91 \\
Qwen 2.5 72B          & 0.93 & 0.76 & 0.93 \\
Qwen 2.5 7B           & 0.72 & 0.72 & 0.16 \\
Llama 3.1 8B          & 0.91 & 0.45 & 0.36 \\
\bottomrule
\end{tabular}
\end{table}

\subsection{Single-policy results}

Table~\ref{tab:main} gives the main results and Table~\ref{tab:recall} recall per decision.

\paragraph{Unsafe approval separates models that accuracy does not.} Qwen 2.5 72B is within five
points of Llama 3.3 70B on accuracy, yet approves twice as many non-coverable cases (18.0\% vs
8.6\%). Llama 3.1 8B approves 46\% of them while recalling 91\% of covered cases: it says
``covered'' to most requests. A leaderboard sorted by accuracy would hide the difference that
matters most to a payer or a patient.

\paragraph{Strong models err toward denial.} The four strongest models keep unsafe approvals at or
below 2.7\%, and most of their remaining errors are over-denials of covered cases (2--17\%). GPT-5
mini has only four, three of them malformed JSON. About half of Gemini 3.1 Flash Lite's 20 are
arithmetic or boundary errors: a patient aged 85 years and 2 months judged over a maximum age of
85, a score of exactly 60\% judged to fail a ``60\% or less'' criterion, ``at least 11 months after
the month of the last test'' counted wrongly, and percentages of a laboratory limit computed
incorrectly.

\paragraph{Case construction showed up as model error.} In our first run, GPT-5 mini over-denied
17 covered cases. Ten asked for documentation of requirements that apply to the program or the
claim rather than the patient (registry participation for TAVR and mitral TEER, claim diagnosis
coding for Pap smears, a post-approval study for blood-based colorectal screening), and three
hospital-bed cases were denied because the requested service, then the NCD's full scope, named
features the prescription did not order. We revised the 49 cases in those five policies to name
the specific service requested and to state that the program- and claim-level requirements are
met, without changing any label, and re-ran them. GPT-5 mini's over-denial fell from 9.8\% to
2.3\% and its accuracy rose from 94.3\% to 96.7\%; the other strong models moved by up to two
points, and unsafe approval rates were essentially unchanged. All results in this paper use the
revised cases. Benchmarks built from policy text should state what the case record is presumed to
establish, or strong models will reasonably ask for it.

\paragraph{Small models cannot use the third label.} Qwen 2.5 7B recovers 16\% of
\insuf{} cases, calling 68 of the 90 not covered. The distinction between ``fails the policy''
and ``does not show that it meets the policy'' appears to require more capacity than the
two-way decision.

\paragraph{The single-policy track is close to saturation for strong models.} 205 of 513 cases
are solved by all eight models, and only two are missed by all four strongest: an
\insuf{} case on renal denervation that all four call \ncov{}, and an insulin pump case at a
C-peptide threshold. Before the revision above, seven cases were missed by all four; five of
those were explained by case construction rather than by the label. We treat the remaining two as
candidate label errors for expert review. The long-context sample, drawn before the revision,
excludes all seven.

\subsection{Long-context results}

\begin{table}[t]
\centering
\small
\caption{Long-context track, 120 cases. ``1 NCD'' is the single-policy result on the same 120
cases. Sizes are estimated from characters; mean actual prompt lengths are about 27--28k and 80k
tokens. $^\dagger$63\% of Llama 3.1 8B's replies at 96k are unparseable, which lowers its unsafe
approval rate mechanically. \tbd{} entries await runs.}
\label{tab:long}
\begin{tabular}{lcccccccc}
\toprule
 & \multicolumn{3}{c}{Accuracy} & \multicolumn{3}{c}{Unsafe appr.\ $\downarrow$} & \multicolumn{2}{c}{NCD acc.} \\
\cmidrule(lr){2-4}\cmidrule(lr){5-7}\cmidrule(lr){8-9}
Model & 1 NCD & 32k & 96k & 1 NCD & 32k & 96k & 32k & 96k \\
\midrule
GPT-5 mini            & 0.958 & \tbd & \tbd & 0.013 & \tbd & \tbd & \tbd & \tbd \\
Gemini 3.1 Flash Lite & 0.942 & 0.900 & \tbd & 0.013 & 0.050 & \tbd & 1.000 & \tbd \\
Qwen 3.7 Plus         & 0.892 & \tbd & \tbd & 0.000 & \tbd & \tbd & \tbd & \tbd \\
DeepSeek V3.1         & 0.883 & \tbd & \tbd & 0.037 & \tbd & \tbd & \tbd & \tbd \\
Llama 3.3 70B         & 0.867 & 0.750 & 0.633 & 0.087 & 0.275 & 0.500 & 0.983 & 0.933 \\
Llama 3.1 8B          & 0.600 & 0.558 & 0.200 & 0.388 & 0.425 & 0.188$^\dagger$ & 0.883 & 0.250 \\
\bottomrule
\end{tabular}
\end{table}

Table~\ref{tab:long} reports the long-context track for the models run so far; the remaining
runs are pending. Three patterns are already clear.

\paragraph{Finding the policy is not the bottleneck for capable models.} Gemini 3.1 Flash Lite names
the governing NCD in all 120 cases at 32k, and Llama 3.3 70B in 98\% at 32k and 93\% at 96k. Yet
both lose accuracy: Gemini from 94.2\% to 90.0\%, Llama 3.3 70B from 86.7\% to 75.0\% and then
63.3\%. Among the cases where Llama 3.3 70B names the right NCD, accuracy is 75.4\% at 32k and
61.6\% at 96k, so most of the loss comes from applying the right policy worse, not from applying
the wrong one.

\paragraph{Long context shifts errors toward approval.} With the policy alone, Llama 3.3 70B's errors
on these 120 cases are split between unsafe approvals (8.7\%) and over-denials (12.5\%). In the
library its unsafe approval rate rises to 27.5\% at 32k and 50\% at 96k, while over-denial falls
to 7.5\% and then zero. Gemini's
rises from 1.3\% to 5.0\% at 32k. Accuracy alone would understate this: the model is not just
noisier, it is more permissive, which is the costly direction for a payer.

\paragraph{Small models break down.} Llama 3.1 8B holds up at 32k (accuracy 55.8\% vs 60.0\% alone)
but collapses at 96k: 63\% of its replies are not valid JSON, it names the right NCD in 25\% of
cases, and accuracy is 20\%. Where it does name the right NCD, it is correct 60\% of the time,
about its single-policy level.

\paragraph{Position.} Llama 3.3 70B at 32k is most accurate when the target is first or last in the
library (22/24 and 19/24) and least accurate in the middle (15/24 at depths 0.5 and 0.75),
consistent with prior findings on position~\citep{liu2024lost}. Gemini shows no such effect (21--22
of 24 at every depth). With 24 cases per depth these differences are suggestive, not conclusive.

\todo{Complete GPT-5 mini, Qwen 3.7 Plus, DeepSeek V3.1 and Gemini at 96k; add confidence intervals.}

\section{Limitations}
\label{sec:limits}

\paragraph{Labels are drafts.} No clinician or certified coder has yet reviewed the labels. They
are derived mechanically from checklists that pass automatic validation, but the checklists
themselves and the reading of each vignette are judgments. We release a review tool and will
report inter-annotator agreement and the share of labels changed. Until then, results should be
read as provisional.

\paragraph{AI-assisted construction.} Cases were drafted with the help of a language model. This
risks favoring models that share its phrasing or assumptions, and may make some cases easier for
models of the same family. We evaluate no model from the drafting family in this version.

\paragraph{Scope.} The benchmark uses NCDs only. Most Medicare coverage in practice is decided
under Local Coverage Determinations, which add documentation and frequency rules; these cannot be
redistributed because of CPT licensing. A model that does well here may still fail on local
policy. Several of our policies have LCD rules that a real reviewer would apply (for example CPAP
adherence thresholds); we record these and keep cases from depending on them.

\paragraph{Presumed facts.} Cases now state that program- and claim-level requirements are met
where a policy has them, but we revised only the five policies where baseline errors exposed the
gap. Other policies may contain requirements a model could reasonably ask to see documented.

\paragraph{Synthetic cases.} Vignettes are short and clean compared with real medical records.
Real prior authorization requires finding facts in long, noisy documentation, which this version
does not test.

\paragraph{Evaluation.} Parse failures are malformed JSON (for example an unescaped quote inside
a string), which we count as wrong rather than repair. Each model saw each case once, through a single router, at temperature 0;
provider-side differences and nondeterminism are not measured. Citation recall uses token overlap
and can credit a quote that shares words with the right span without being it.

\paragraph{Contamination.} All data are public. We plan a held-out test set with private labels.

\section{Ethical considerations}

MedCoverBench measures reasoning over policy text. It is not a coverage determination tool and
should not be used to decide anyone's coverage. We report unsafe approval as the headline metric
because our concern is systems that approve what policy denies, but over-denial harms patients as
well, and we report both. The cases are synthetic and contain no protected health information.

\section{Conclusion}

MedCoverBench tests whether a language model can do the core step of coverage review: apply a
real written policy to a case, cite the rule, and recognize missing documentation. On single
policies, strong models are close to the ceiling while small open models approve many cases the
policy denies; unsafe approval exposes differences that accuracy hides. The long-context track,
where the policy must first be found among many others, shows that in the models run so far
the policy is usually found but then applied more permissively, so unsafe approvals rise sharply
with context length even when retrieval succeeds. The next steps are expert
review of every label, a held-out test set, and harder cases that combine several criteria and
exceptions.

\paragraph{Availability.} Code, data, model outputs and the build log:
\url{https://github.com/mdobbali/medcoverbench}.

\bibliographystyle{plainnat}
\bibliography{refs}

\appendix

\section{Prompt}
\label{app:prompt}
\small
\begin{verbatim}
You are a Medicare coverage reviewer. Decide whether the requested service
is covered under the National Coverage Determination (NCD) provided, using
only that policy and the case facts.

Decisions:
- "covered": every applicable coverage criterion is met and documented.
- "not_covered": the facts show at least one applicable criterion is not met.
- "insufficient_documentation": no criterion is shown to be violated, but
  documentation required by the policy is missing.

Respond with only a JSON object:
{"decision": "covered" | "not_covered" | "insufficient_documentation",
 "policy_quotes": ["verbatim text from the policy for each criterion that
                   determines the decision"],
 "missing_documentation": ["each required item absent from the case"],
 "rationale": "2-4 sentences"}
\end{verbatim}
The user message contains the NCD (title, then each field under a heading), the requested service
and the vignette. The long-context prompt is the same, except that it describes a library of NCDs,
asks the model first to find the governing NCD, and adds an \texttt{ncd} field to the output.
\normalsize

\section{Policies}
\label{app:policies}
\small
\begin{longtable}{llrrrr}
\toprule
NCD & Title & Criteria & Cov. & Not cov. & Insuf. \\
\midrule
\endhead
20.4 & Implantable Cardioverter Defibrillators (ICDs) & 13 & 4 & 5 & 1 \\
20.6 & Transmyocardial Revascularization (TMR) & 9 & 3 & 5 & 2 \\
20.7 & Percutaneous Transluminal Angioplasty (PTA) & 9 & 4 & 5 & 2 \\
20.9.1 & Ventricular Assist Devices & 9 & 4 & 5 & 2 \\
20.29 & Hyperbaric Oxygen Therapy & 9 & 4 & 5 & 2 \\
20.32 & Transcatheter Aortic Valve Replacement (TAVR) & 10 & 3 & 6 & 1 \\
20.33 & Transcatheter Edge-to-Edge Repair (TEER) for Mitral Valve\ldots & 11 & 4 & 5 & 1 \\
20.34 & Percutaneous Left Atrial Appendage Closure (LAAC) & 9 & 3 & 6 & 1 \\
20.35 & Supervised Exercise Therapy (SET) for Symptomatic\ldots & 8 & 3 & 5 & 2 \\
20.36 & Implantable Pulmonary Artery Pressure Sensors for Heart\ldots & 10 & 3 & 5 & 2 \\
20.37 & Transcatheter Tricuspid Valve Replacement (TTVR) & 6 & 3 & 5 & 2 \\
20.38 & Transcatheter Edge-to-Edge Repair for Tricuspid Valve\ldots & 6 & 3 & 5 & 2 \\
20.39 & Cardiac Contractility Modulation (CCM) for Heart Failure (HF) & 6 & 3 & 5 & 2 \\
20.40 & Renal Denervation (RDN) for Uncontrolled Hypertension & 9 & 2 & 7 & 1 \\
50.3 & Cochlear Implantation & 5 & 3 & 5 & 2 \\
70.2.1 & Services Provided for the Diagnosis and Treatment of\ldots & 6 & 3 & 4 & 2 \\
90.1 & Pharmacogenomic Testing for Warfarin Response & 7 & 3 & 5 & 2 \\
90.2 & Next Generation Sequencing (NGS) & 11 & 3 & 5 & 2 \\
100.1 & Bariatric Surgery for Treatment of Co-Morbid Conditions\ldots & 9 & 3 & 5 & 2 \\
110.4 & Extracorporeal Photopheresis & 4 & 3 & 5 & 2 \\
110.14 & Apheresis (Therapeutic Pheresis) & 9 & 4 & 5 & 2 \\
110.23 & Stem Cell Transplantation (Formerly 110.8.1) & 14 & 4 & 7 & 2 \\
150.2 & Osteogenic Stimulators & 9 & 4 & 6 & 2 \\
160.7 & Electrical Nerve Stimulators & 6 & 3 & 5 & 2 \\
160.12 & Neuromuscular Electrical Stimulation (NMES) & 9 & 3 & 6 & 2 \\
160.18 & Vagus Nerve Stimulation (VNS) & 9 & 4 & 5 & 2 \\
160.24 & Deep Brain Stimulation for Essential Tremor and\ldots & 15 & 4 & 6 & 3 \\
190.14 & Human Immunodeficiency Virus (HIV) Testing (Diagnosis) & 9 & 4 & 4 & 1 \\
200.3 & Monoclonal Antibodies Directed Against Amyloid for the\ldots & 6 & 4 & 5 & 1 \\
210.2 & Screening Pap Smears and Pelvic Examinations for Early\ldots & 8 & 4 & 3 & 2 \\
210.3 & Colorectal Cancer Screening Tests & 13 & 3 & 5 & 2 \\
210.7 & Screening for the Human Immunodeficiency Virus (HIV) Infection & 7 & 4 & 4 & 2 \\
210.12 & Intensive Behavioral Therapy for Obesity & 8 & 3 & 5 & 2 \\
210.13 & Screening for Hepatitis C Virus (HCV) in Adults & 6 & 3 & 5 & 2 \\
210.14 & Lung Cancer Screening with Low Dose Computed Tomography (LDCT) & 8 & 3 & 4 & 1 \\
220.6.1 & PET for Perfusion of the Heart & 7 & 4 & 4 & 2 \\
220.6.13 & FDG PET for Dementia and Neurodegenerative Diseases & 11 & 3 & 5 & 2 \\
230.10 & Incontinence Control Devices & 8 & 3 & 5 & 2 \\
240.1 & Lung Volume Reduction Surgery (Reduction Pneumoplasty) & 12 & 3 & 5 & 2 \\
240.2 & Home Use of Oxygen & 9 & 4 & 4 & 2 \\
240.4 & Continuous Positive Airway Pressure (CPAP) Therapy For\ldots & 9 & 4 & 5 & 1 \\
240.9 & Noninvasive Positive Pressure Ventilation (NIPPV) in the\ldots & 7 & 5 & 5 & 0 \\
260.1 & Adult Liver Transplantation & 8 & 4 & 5 & 2 \\
260.3 & Pancreas Transplants & 7 & 4 & 5 & 2 \\
260.9 & Heart Transplants & 5 & 4 & 4 & 2 \\
280.3 & Mobility Assistive Equipment (MAE) & 10 & 4 & 5 & 2 \\
280.6 & Pneumatic Compression Devices & 8 & 3 & 5 & 2 \\
280.7 & Hospital Beds & 8 & 4 & 4 & 2 \\
280.8 & Air-Fluidized Bed & 11 & 3 & 5 & 2 \\
280.14 & Infusion Pumps & 12 & 4 & 5 & 2 \\
\bottomrule
\end{longtable}
\normalsize

\end{document}
